\section{Modèle de données et diagramme de classes}

\subsection{Principes de modélisation}

Le modèle conceptuel est organisé autour de quatre notions centrales :
l'organisation qui exploite le réseau, la station qui fournit le service, le
client qui utilise ce service et la session qui matérialise la recharge. Les autres
classes complètent ce noyau pour couvrir l'exploitation et les deux relations
commerciales de la plateforme.

La notation suit les conventions UML usuelles : une classe est séparée en
compartiments, les attributs portent un nom et un type, et chaque extrémité
d'association indique sa multiplicité \cite{audibertUmlClasses}. Une composition
n'est employée que lorsque la partie dépend du cycle de vie du tout. Les attributs
précédés de \texttt{/} sont dérivés de preuves techniques plutôt que saisis
directement.

La conception respecte les principes de normalisation relationnelle. Une
information métier n'est portée que par l'entité qui en est responsable : les
connecteurs sont séparés des stations, les paiements des sessions et les offres des
adhésions. Les relations plusieurs-à-plusieurs sont résolues par des entités
d'association, notamment l'adhésion d'un client à un plan et l'abonnement d'une
organisation à une offre SaaS. Cette organisation évite les groupes répétitifs et
les dépendances partielles.

Le diagramme reste volontairement conceptuel et ne contient pas les opérations de
framework. Les tables de journalisation, les
événements OCPP, les échantillons de mesure, les fichiers et les notifications ne
sont pas représentés afin de préserver sa lisibilité. Ces éléments techniques
conservent l'historique, mais ne modifient pas la structure des concepts métier
principaux.

\subsection{Domaines fonctionnels}

\subsubsection{Identité et séparation organisationnelle}

Une organisation possède ses employés et son patrimoine de recharge. Les
administrateurs, opérateurs et techniciens sont rattachés à une seule organisation.
Le super administrateur et le client sont des acteurs globaux. Un même client peut
utiliser les stations de plusieurs organisations ; son compte n'est donc pas
rattaché à l'une d'elles. La relation avec une organisation apparaît au travers des
adhésions et des sessions.

\subsubsection{Réseau et recharge}

Une organisation exploite plusieurs stations. Chaque station possède au moins un
connecteur identifié dans le protocole OCPP. Une session associe un client à un
connecteur pendant une période déterminée et conserve l'énergie délivrée, le coût
calculé et son état. Elle référence le tarif appliqué et, lorsqu'il existe,
l'avantage provenant d'une adhésion. Une session finalisée possède au plus un
paiement principal.

L'état de disponibilité n'est pas une simple saisie. Il est projeté à partir des
messages OCPP récents, des connecteurs, des opérations de maintenance et des
éventuelles dérogations. Les traces techniques sont conservées séparément pour
expliquer chaque transition.

\subsubsection{Exploitation}

Une station peut produire plusieurs alertes et être couverte par des plans de
maintenance. Une alerte ou une échéance de maintenance peut conduire à une
intervention. Celle-ci porte l'affectation du technicien, la priorité, le délai SLA
et l'état de traitement. Les événements, photos, documents et rapports
d'intervention constituent des preuves associées sans alourdir le diagramme global.

\subsubsection{Commerce client et commerce SaaS}

Deux cycles commerciaux sont séparés. Le premier relie le client à une organisation
au moyen d'un plan de recharge et d'une adhésion ; il peut accorder une remise sur
les sessions. Le second relie l'organisation à la plateforme au moyen d'une offre
SaaS et d'un abonnement commercial ; il détermine la période d'essai, les capacités
et le droit d'utiliser les fonctions de gestion. Cette distinction empêche de
confondre le paiement d'une recharge et la facturation du logiciel.

\subsection{Contraintes et invariants principaux}

\begin{longtable}{|p{3.7cm}|p{10.8cm}|}
  \caption{Invariants structurants du modèle}
  \label{tab:domain-invariants}\\
  \hline
  \textbf{Domaine} & \textbf{Invariant}\\
  \hline
  \endfirsthead
  \hline
  \textbf{Domaine} & \textbf{Invariant}\\
  \hline
  \endhead
  Identité &
  Un utilisateur actif possède exactement un rôle ; un rôle organisationnel exige
  une organisation active, alors qu'un rôle global interdit cette association.\\\hline
  Multi-tenant &
  Toute lecture ou écriture opérationnelle est limitée à l'organisation de
  l'utilisateur, sauf permission explicite du super administrateur.\\\hline
  Commissioning &
  Une station et ses connecteurs sont créés dans une même opération ; la référence
  et l'identité OCPP sont uniques.\\\hline
  Disponibilité &
  Le statut métier est dérivé des preuves OCPP récentes et des décisions
  opérationnelles ; il n'est pas saisi directement pour une station connectée.\\\hline
  Recharge &
  Un client et un connecteur ne peuvent participer qu'à une tentative ou une
  session active à la fois.\\\hline
  Mesure &
  L'énergie cumulée d'une session ne diminue jamais et chaque valeur reste
  rattachée à la transaction OCPP correspondante.\\\hline
  Paiement &
  Une même demande de paiement ne peut être traitée qu'une fois et un événement du
  fournisseur ne modifie qu'une transaction connue.\\\hline
  Tarification &
  Le tarif et la remise sont figés au début de la session afin que son coût reste
  explicable.\\\hline
  Intervention &
  Le technicien affecté, la station et l'intervention appartiennent à la même
  organisation.\\\hline
  Abonnement &
  L'abonnement client et l'abonnement SaaS suivent des cycles indépendants et ne
  partagent aucune facture.\\
  \hline
\end{longtable}

\clearpage
\subsection{Diagramme de classes global}

La figure~\ref{fig:conceptual-domain-model} présente une vue unique des classes
métier principales. Les attributs sont limités aux informations nécessaires pour
comprendre la responsabilité de chaque classe. Les relations utilisent des tracés
orthogonaux ou directs, sans boucle réflexive ni circuit décoratif. Les cardinalités
décrivent les relations fonctionnelles et non les détails propres à Eloquent ou aux
migrations.

\begin{figure}[H]
  \centering
  \resizebox{0.96\textwidth}{!}{\begin{tikzpicture}[
  font=\sffamily\tiny,
  class/.style={
    draw=ctline, rounded corners=1.5mm, thick, fill=white,
    text width=29mm, minimum height=16mm, align=left, inner sep=2mm
  },
  identity/.style={class, draw=ctblue, fill=blue!5},
  network/.style={class, draw=ctgreen, fill=ctmint},
  operations/.style={class, draw=ctblue, fill=cyan!4},
  commerce/.style={class, draw=ctamber, fill=yellow!7},
  rel/.style={thick, draw=ctgray},
  card/.style={fill=white, inner sep=1pt, font=\sffamily\tiny, text=ctgray}
]
  \node[identity] (user) at (-4.8,10.0)
    {\textbf{Utilisateur}\\\rule{28mm}{0.2pt}\\identité, rôle\\statut};
  \node[identity] (org) at (0,10.0)
    {\textbf{Organisation}\\\rule{28mm}{0.2pt}\\nom, statut\\coordonnées};
  \node[commerce] (saas) at (4.8,10.0)
    {\textbf{Offre SaaS}\\\rule{28mm}{0.2pt}\\capacités, prix\\statut};

  \node[commerce] (plan) at (-4.8,7.0)
    {\textbf{Plan de recharge}\\\rule{28mm}{0.2pt}\\abonnement, remise\\validité};
  \node[network] (station) at (0,7.0)
    {\textbf{Station}\\\rule{28mm}{0.2pt}\\référence, position\\disponibilité};
  \node[commerce] (orgsub) at (4.8,7.0)
    {\textbf{Abonnement organisation}\\\rule{28mm}{0.2pt}\\essai, période\\état};

  \node[commerce] (plansub) at (-4.8,4.0)
    {\textbf{Adhésion client}\\\rule{28mm}{0.2pt}\\période, renouvellement\\état};
  \node[network] (connector) at (0,4.0)
    {\textbf{Connecteur}\\\rule{28mm}{0.2pt}\\identifiant OCPP\\type, puissance, état};
  \node[operations] (maintenance) at (4.8,4.0)
    {\textbf{Plan de maintenance}\\\rule{28mm}{0.2pt}\\fréquence, échéance\\état};

  \node[commerce] (tariff) at (-4.8,1.0)
    {\textbf{Tarif}\\\rule{28mm}{0.2pt}\\prix énergie, frais\\période de validité};
  \node[network] (session) at (-0.4,1.0)
    {\textbf{Session de recharge}\\\rule{28mm}{0.2pt}\\début, fin, énergie\\coût, état};
  \node[operations] (alert) at (3.5,1.0)
    {\textbf{Alerte}\\\rule{28mm}{0.2pt}\\type, sévérité\\état};

  \node[commerce] (payment) at (-0.4,-2.0)
    {\textbf{Paiement}\\\rule{28mm}{0.2pt}\\référence, montant\\état};
  \node[operations] (intervention) at (4.8,-2.0)
    {\textbf{Intervention}\\\rule{28mm}{0.2pt}\\priorité, SLA\\affectation, état};

  \draw[rel] (user) -- node[card, above]{0..* employés \quad 0..1} (org);
  \draw[rel] (org) -- node[card, right]{1 \quad 0..*} (station);
  \draw[rel] (org.south west) -- node[card, above left]{1 \quad 0..*} (plan.north east);
  \draw[rel] (org.south east) -- node[card, above right]{1 \quad 0..1} (orgsub.north west);
  \draw[rel] (saas) -- node[card, right]{1 \quad 0..*} (orgsub);

  \draw[rel] (plan) -- node[card, right]{1 \quad 0..*} (plansub);
  \draw[rel] (user.west) -- ++(-0.55,0) |- node[card, near end, left]
    {1 \quad 0..*} (plansub.west);

  \draw[rel] (station) -- node[card, right]{1 \quad 1..*} (connector);
  \draw[rel] (station.south east) -- node[card, above right]{1 \quad 0..*}
    (maintenance.north west);
  \draw[rel] (station.south east) to[out=-35,in=105]
    node[card, above]{1 \quad 0..*} (alert.north);

  \draw[rel] (connector) -- node[card, right]{1 \quad 0..*} (session);
  \draw[rel] (user.south west) -- ++(-0.9,0) |- (-6.2,-0.15) --
    node[card, below]{client 1 \quad 0..*} (session.south west);
  \draw[rel] (plansub.south east) to[out=-25,in=155]
    node[card, above]{0..1 \quad 0..*} (session.north west);
  \draw[rel] (tariff) -- node[card, above]{0..1 \quad 0..*} (session);
  \draw[rel] (session) -- node[card, right]{1 \quad 0..1} (payment);

  \draw[rel] (alert) -- node[card, above right]{0..1 \quad 0..*}
    (intervention.north west);
  \draw[rel] (maintenance.east) -- ++(0.45,0) |-
    node[card, near end, right]{0..1 \quad 0..*} (intervention.east);
\end{tikzpicture}
}
  \caption{Diagramme de classes conceptuel global de \ProjectName{}}
  \label{fig:conceptual-domain-model}
\end{figure}
